\subsection{Structural model parameterizations}\label{sec:parameterization}
This section describes the functional form assumptions we make in estimating the structural model.\par 
The production functions of the PSU score and of GPA are as follows:
\begin{equation}
    PSU_i=\beta_0^P+\beta_1^P e_i +\beta_2^P y_{i,t-1}^{(1)}+\epsilon_{i}^P,
\end{equation}
\begin{equation}
    GPA_i=\beta_0^G+\beta_1^g e_i +\beta_2^G y_{i,t-1}^{(2)}+\epsilon_{i}^G,
\end{equation}
\noindent where $y_{i,t-1}^{(1)}$ is a baseline standardized test score and $y_{i,t-1}^{(2)}$ is the baseline GPA (we restrict $GPA_i$ to be between 1 and 7). We assume that the technology shocks $\epsilon_i=[ \epsilon_i^P,\epsilon_i^G]$ are distributed as bivariate normal: $\epsilon_{it}\sim N(0,\Sigma)$, with $\Sigma=\begin{bmatrix}
	\sigma^{2}_{P} & \rho\sigma_{P}\sigma_{G}   \\
	\rho\sigma_{P}\sigma_{G}  & \sigma^{2}_{G} \\
\end{bmatrix}$. Given a PSU score, the probability of a regular admission is 
\begin{equation}\label{eq:regular_admission}
    Pr(A^R_i=1|PSU_i,S_i=1;\gamma)= \Phi(\gamma_0 + \gamma_1 PSU_i).
\end{equation}

The subjective production functions of the PSU score and of GPA are as follows:
\begin{equation}\label{eq:rescalePSUb}
	PSU^b_{it}=\beta^{Pb}_{0k_i}+\beta^{Pb}_{1k_i} e_{it} +\beta^{Pb}_2 y^{(1)}_{it-1}+\epsilon^{PSU^b}_{it}, \quad \epsilon_{it}^{PSU^b}\sim N(0,\sigma^{2}_{PSU^b})
\end{equation}
\begin{equation}\label{eq:rescaleGPAb}
	GPA^b_{it}=\beta^{Gb}_{0}+\beta^{Gb}_{1k_i} e_{it} +\beta^{Gb}_2 y^{(2)}_{it-1}+\epsilon^{GPA^b}_{it}, \quad \epsilon_{it}^{GPA^b}\sim N(0,\sigma^{2}_{GPA^b})
\end{equation}

\noindent where the shocks $(\epsilon^{PSU^b}_{it},\epsilon^{GPA^b}_{it}) $ are i.i.d. normal and capture belief uncertainty. Observationally identical students hold heterogeneous beliefs about the production function: parameters $\beta^{Pb}_{0k_i},\beta^{Pb}_{1k_i}, \beta^{Gb}_{1k_i}$ vary with the student's unobserved type. The believed outcomes vary also with baseline characteristics and effort.\par 

The subjective probability of regular admission, conditional on taking the PSU entrance exam ($S_i=1$), is equal to the subjective probability that a student's believed score will be above the believed admission cut-off point. Students form a subjective probability distribution for the admission cutoff: $c_i^{Rb}\sim N(\bar{c}^{Rb},\sigma^2_{c^{Rb}})$. Letting $\overline{PSU}^b_{it}=\beta^{Pb}_{0k_i}+\beta^{Pb}_{1k_i} e_{it} +\beta^{Pb}_2 y^{(1)}_{it-1}$ denote the expected PSU score, $\epsilon_i^{c^{Rb}}$ the mean-zero additive belief shock around the expected cutoff, and $A^R_i$ a dummy for a regular admission, the subjective probability of a regular admission is:
\vspace{-12pt}
\begin{eqnarray}\label{eq:pr_Rb}
	Pr^b(A^R_i=1|e_{it},y_{it-1}^{(1)},k_i,S_i=1)&=&Pr\left( \overline{PSU}^b_{it}+\epsilon_i^{PSU^b}\geq \bar{c}^{Rb}+\epsilon_i^{c^{Rb}} \right)\\ \nonumber
	&=&\Phi \left( \frac{\overline{PSU}^b_{it}-\bar{c}^{Rb}}{\sqrt{\sigma^{2}_{PSU^b}+\sigma^2_{c^{Rb}}}}\right) \\ \nonumber
	&=&\Phi\left(\gamma_0^b +\gamma_1^b \overline{PSU}^b_{it}\right),
\end{eqnarray}
\noindent where $\gamma_0^b=\frac{-\bar{c}^{Rb}}{\sqrt{\sigma^{2}_{PSU^b}+\sigma^2_{c^{Rb}}}}$ and $\gamma_1^b=\frac{1}{\sqrt{\sigma^{2}_{PSU^b}+\sigma^2_{c^{Rb}}}}$ and $\Phi(\cdot)$ is the standard Normal cumulative distribution function. Given an expected PSU score, uncertainty is generated by uncertainty around own score ($\sigma^{2}_{PSU^b}$) and around the admission cutoff ($\sigma^2_{c^{Rb}}$), which are absorbed by the parameters $\gamma_0^b$ and $\gamma_1^b$. As it is standard  to impose functional form restrictions on subjective probabilities (e.g. \citealp{delavande2019university,kapor2020heterogeneous}), we impose normality.\par 

Letting $\overline{GPA}^b_{it}=\beta^{Gb}_{0}+\beta^{Gb}_{1k_i} e_{it} +\beta^{Gb}_2 y^{(2)}_{it-1}$ denote the expected GPA, $\epsilon_i^{c15b}$ the mean-zero belief shock around the expected school cutoff\footnote{Students form a subjective probability distribution for the cutoff in their school:  $c15^b_i\sim N(\bar{c15}^b_i,\sigma^2_{c15^b})$, characterized by a heterogeneous expected cutoff, $\bar{c15}^b_i$, with uncertainty around it, $\sigma^2_{c15^b}$. We assume our survey instrument measured the expected cutoff $\bar{c15}^b_i$ for each student $i$. The elicited $\bar{c15}^b_i$ is missing for less than $20\%$ of students. We assume these students correctly predict the cutoff; thus, results provide a lower bound to the role that biased rank beliefs play in policy response.}, and $A_i^P$ a dummy for a preferential admission, the subjective probability of a preferential admission, conditional on taking the entrance exam ($S_i=1$), for students in treated schools is:
\vspace{-10pt}
\begin{eqnarray}\label{eq:pr_Pb}
	Pr^b(A^P_i=1|e_{it},y_{it-1}^{(2)},k_i,S_i=1)&=&Pr\left( \overline{GPA}^b_{it}+\epsilon_i^{GPA^b}\geq c_0+\bar{c15}^b_i+\epsilon_i^{c15b} \right)\\ \nonumber
	&=&\Phi \left( \frac{\overline{GPA}^b_{it}-c_0- \bar{c15}_i^b}{\sqrt{\sigma^2_{GPA^b}+\sigma^2_{c15^b}}}\right) \\ \nonumber
	&=&\Phi\left(\xi_0^b+\xi_1^b (\overline{GPA}^b_{it}-\bar{c15}_i^b)\right),
\end{eqnarray}
\noindent where $\xi_0^b=\frac{-c_0}{\sqrt{\sigma^2_{GPA^b}+\sigma^2_{c15^b}}}$ and $\xi_1^b=\frac{1}{\sqrt{\sigma^2_{GPA^b}+\sigma^2_{c15^b}}}$.\footnote{Parameter $c_0$ is a net adjustment to the GPA and the cutoff to capture the fact that the top $15\%$ rule is based on adjusted GPA. } Given an expected $GPA$ and an expected cutoff, uncertainty is generated by the uncertainty around own GPA ($\sigma^2_{GPA^b}$) and around the school cutoff ($\sigma^2_{c15^b}$), which are absorbed by the parameters $\xi_0^b$ and $\xi_1^b$. As before, we assume normality. \par 

In the first period, the per-period utility from effort depends on how effort affects achievement. We assume achievement is produced as follows: $y_{i}=\alpha_{0k_i}+ \alpha_1 x_i +\alpha_2 e_{it} +\alpha_3$. We assume that our survey measures study effort with additive noise:  $e_i^o=e_i+\epsilon_i^{m.e.e.}$, where $\epsilon^{m.e.e.}\sim N(0,\sigma^2_{m.e.e.})$ is a classical measurement error. We assume that our standardized test score measures achievement with additive noise: $y_i^o=y_i+\epsilon_i^{m.e.y.}$, with $\epsilon_i^{m.e.y.}\sim N(0,\sigma^2_{m.e.y.})$.\par 

As in the real-world admission system, the selectivity of an admission depends on a student's PSU (for regular admissions) and GPA (for preferential admissions). We assume the following functional forms:
\begin{equation}\label{eq:selR}
    q^R(PSU_i)=\lambda^R_0 +\lambda^R_1 PSU_i + \epsilon_i^{qR}
\end{equation}
\begin{equation}\label{eq:selP}
    q^P(GPA_i)=\lambda^P_0 +\lambda^P_1 GPA_i + \epsilon_i^{qP}.
\end{equation}
  

